\documentclass[11pt]{article}
\usepackage[T1]{fontenc}
\usepackage{lmodern,amsmath,amssymb,amsthm,booktabs,microtype}
\usepackage[a4paper,margin=26mm]{geometry}
\usepackage[hidelinks]{hyperref}
\newtheorem{theorem}{Theorem}[section]
\newtheorem{proposition}[theorem]{Proposition}
\newtheorem{lemma}[theorem]{Lemma}
\newtheorem{corollary}[theorem]{Corollary}
\newcommand{\Q}{\mathbb Q}
\newcommand{\Z}{\mathbb Z}
\newcommand{\C}{\mathbb C}
\newcommand{\R}{\mathbb R}
\newcommand{\OO}{\mathcal O}
\newcommand{\tr}{\operatorname{tr}}
\newcommand{\Frob}{\operatorname{Frob}}
\newcommand{\ord}{\operatorname{ord}}
\newcommand{\workbenchurl}{https://github.com/KokunoYumeto/zeta-and-l-functions-workbench}
\newcommand{\recordurl}{https://github.com/KokunoYumeto/mathematics-workbenches}
\title{Real nonvanishing and prime-sum tests for a cubic Hecke $L$-function, after Louboutin and Munsch}
\author{}
\date{24 November 2026}
\begin{document}
\maketitle
\begin{abstract}
Louboutin and Munsch gave a finite partial-sum criterion that makes a quadratic theta function positive and excludes positive real zeros of its Dirichlet $L$-function. We apply the same idea to the weight-one theta series of the pure cubic field generated by the cube root of 269. Exact agreement of two coefficient computations, followed by a rational tail bound, proves positivity of the theta kernel and of the completed cubic Hecke $L$-function on the whole real axis. In particular its central value is nonzero. We also establish the precise signed-square, prime-window and squarefree-composite tests for its Riemann hypothesis, and quantify why the theta positivity and ordinary Rankin--Selberg positivity do not supply their missing bound. The mathematics and text were produced by GPT-6 Astra in ChatGPT Pro mode; the checks listed in the workbench are by the same model, and no human has checked this text.
\end{abstract}

\section{Main statements and the unresolved Riemann hypothesis}
This section identifies the function, the unconditional conclusion, and the stronger estimate that remains open.
Let
\[
 F=\Q(\sqrt[3]{269}),\qquad K=\Q(\omega),\qquad E=FK,
 \qquad \omega^2+\omega+1=0.
\]
The splitting group is $S_3$. Its standard two-dimensional representation $\rho$ is the sum-zero part of the action on the three roots. Write
\[
 L(s)=L(s,\rho)=\frac{\zeta_F(s)}{\zeta(s)},\quad
 Q=3\cdot269^2=217083,
 \quad\Lambda(s)=2\left(\frac{\sqrt Q}{2\pi}\right)^s\Gamma(s)L(s).
 \tag{1.1}
\]
The quotient is entire because $\rho$ is induced from a nontrivial cubic Hecke character of $K$. The completed function is entire and satisfies $\Lambda(s)=\Lambda(1-s)$. These are classical Hecke theory and Artin formalism; no Riemann hypothesis is assumed. A full author-hosted proof of the analytic input is Poonen (2015, Theorem 5.22(b),(c), pp.~38--39), using his global zeta-integral Theorem 5.16.

\begin{theorem}\label{thm:positive}
Write $L(s)=\sum_{n\ge1}A(n)n^{-s}$ in its half-plane of absolute convergence and set
\[
 f(t)=\sum_{n\ge1}A(n)\exp\left(-\frac{2\pi nt}{\sqrt Q}\right).
\]
Then
\[
 f(t)\ge\frac{999999}{1000000}
            \exp\left(-\frac{2\pi t}{\sqrt Q}\right)>0
 \qquad(t\ge1).
 \tag{1.2}
\]
For every real $s$, $\Lambda(s)>0$. Consequently $L(s)>0$ for $s>0$, and $L$ has no real zero in its critical strip. A rigorous enclosure is
\[
 \left|L(1/2)-3.585687918544579072941074961432534305815704362\right|
       < 9.3\cdot10^{-46}.
 \tag{1.3}
\]
\end{theorem}

The finite computation in this theorem is part of its proof, with exact integer assertions that can be rerun. The two implementations use, respectively, rational prime factorization and enumeration in the Eisenstein lattice. The interval in (1.3) is not needed for nonvanishing.

Louboutin and Munsch (2013, Theorem 4, pp.~646--647) prove the corresponding finite-partial-sum implication for theta series of real even Dirichlet characters. Their kernel has exponents $n^2$ and coefficients bounded by one. The proof below gives the variant with exponents $n$, coefficients bounded by the divisor function, and the conductor $Q$ in (1.1). Thus the positivity mechanism is theirs; the exact cubic coefficient and tail certificate is the calculation made here. We make no claim of priority for this specialization.

OpenAI's quasi-Riemann theorem (2026, Theorem 1.1, p.~4) excludes zeros with real part greater than $7/8$ for every finite-order Hecke character of $K$, including this one. It is used below as a named external theorem. Liu's later main theorem (2026, Theorems~1.2 and 8.29, pp.~3 and 78--80) gives the sharper boundary $B_*=0.873345856106718576801184796065\ldots$ for every unitary Hecke character over every number field. His equations (42) and (44), pp.~27--29, define this degree-nine algebraic constant; the indicated root satisfies
\[
 0.8733458561067185768<B_*<0.8733458561067185769.
\]
Liu's theorem is also granted as an external input. Our real positivity conclusion treats all remaining real arguments as well. The distinction matters for criteria whose only possible obstruction is a real value such as $2/3$; the ideal-shell criterion in Section~8 is one example.

The Riemann hypothesis for $L$ remains stronger: it asks that every nontrivial \emph{complex} zero have real part $1/2$. Real nonvanishing eliminates all real candidates, including a possible central zero, but does not control the remaining nonreal zeros. The prime statistic that detects those zeros is particularly simple. For good primes put
\[
 t(p)=\#\{a\bmod p:a^3=269\}-1.
 \tag{1.4}
\]
It equals $2$ or $-1$ at $p\equiv1\pmod3$, according as 269 is or is not a cube, and equals zero at $p\equiv2\pmod3$. Good means $p\nmid3\cdot269m$ when a modulus $m$ is in use.


\subsection*{Signed-square and prime-window criteria}
Fix a positive integer $m$ divisible by 3. Every character
$\chi\bmod m$ is interpreted through its complete primitive Artin
factor, including the ramified local factors. Define
\[
 \begin{gathered}
 J_m(s)=\prod_{\chi\bmod m}L(s,\rho\otimes\chi),\\
 R_m(x)=\sum_{\substack{p\le x,\ p\nmid3\cdot269m\\
                         p\equiv1\bmod m}}t(p)\log p,
 \\
 T_m(K)=K^{-2}\sum_{k\le K}R_m(k^2)
 \end{gathered}
\]
for positive integers $K$. Put $B=\Q(\zeta_m)$ and $T=EB$.
Every factor of $J_m$ is an entire nontrivial Hecke function, and
$J_m=\zeta_T/\zeta_B$.
\begin{theorem}\label{thm:squares}
The following two assertions are equivalent:
\begin{enumerate}
\item Every nontrivial zero of every $L(s,\rho\otimes\chi)$ with $\chi\bmod m$ has real part $1/2$.
\item $\liminf_{K\to\infty}T_m(K)>-\infty$.
\end{enumerate}
No Dirichlet Riemann hypothesis is needed for this relative equivalence. Full Dedekind RH for $T$ also requires RH for the abelian factors of $\zeta_B$.
\end{theorem}

Here are the other main equivalent tests, with all parameters fixed
before a window is chosen. Fix $\delta>0$, put
$P_m(y)=R_m(2y)-R_m(y)$, and let $v_-=\max(0,-v)$.
Write
\[
 -\frac{J_m'}{J_m}(s)=\sum_{n\ge1}\frac{b_m(n)}{n^s}
           \quad(\Re s>1).
\]
Let $b_{m,1}$ agree with $b_m$ at first powers of good primes and
vanish elsewhere, and put $h_m=b_m-b_{m,1}$. Define the centering
constant and the prime ramp by
\[
 \begin{gathered}
 C_m=\frac{J_m'}{J_m}(1)+\sum_{n\ge1}\frac{h_m(n)}n,\\
 V_m(y)=\sum_{p\le y}b_{m,1}(p)(y/p-1)+yC_m.
 \end{gathered}
\]
The series defining $C_m$ converges absolutely. The separate
notation $C_m(y)$ below denotes a composite sum, not this constant.
For $z=(\log y)^2$, let
\[
 C_m(y)=\sum_{\substack{y<n\le2y,\ n\text{ squarefree}\\
             \omega(n)\ge2,\ n\equiv1\bmod m\\
             p\mid n\Rightarrow p>z,\ p\nmid3\cdot269m}}
                     \left(\prod_{p\mid n}t(p)\right)\log n.
\]
Here $\omega(n)$ counts distinct prime factors.
For every fixed $m$ and $\epsilon>0$, one has unconditionally
\[
 P_m(y)+C_m(y)=O_{m,\epsilon}(y^{1/2+\epsilon}).
\]
For these fixed $m$ and $\delta$, each of the following assertions
is equivalent to the twist RH in Theorem~\ref{thm:squares}:
\begin{enumerate}
\item For every $0<\epsilon<1/2$, some constant
$C_{m,\epsilon,\delta}$ and arbitrarily large $Y$ satisfy
\[
 P_m(y)\ge-C_{m,\epsilon,\delta}y^{1/2+\epsilon}
       \qquad(Y\le y\le Y^{1+\delta}).
\]
\item For every $\epsilon>0$, arbitrarily large integers $K$
satisfy
\[
 \sum_{K\le k\le K^{1+\delta}}\frac{(T_m(k))_-}{k}\le K^\epsilon.
\]
\item For every $0<\epsilon<1/2$, some constant
$C_{m,\epsilon,\delta}$ and arbitrarily large $Y$ satisfy
\[
 V_m(y)\le C_{m,\epsilon,\delta}y^{1/2+\epsilon}
       \qquad(Y\le y\le Y^{1+\delta}).
\]
\item For every $0<\epsilon<1/2$, some constant
$C_{m,\epsilon,\delta}$ and arbitrarily large $Y$ satisfy
\[
 C_m(y)\le C_{m,\epsilon,\delta}y^{1/2+\epsilon}
       \qquad(Y\le y\le Y^{1+\delta}).
\]
\end{enumerate}
The window equivalences use Alp\"oge's general Landau-window
theorem \cite{AlpogeWindow}; the arithmetic reductions and their
prime-power and sampling errors are proved in Sections~4--6.
No uniformity in $m$ is asserted. At $m=3$, $J_3=L^2$, so these
are criteria for the one cubic RH. The composite sum includes
every admissible number of prime factors. Replacing the ramp's
centering constant by a value with error $\eta$ changes it by
$\eta y$; a test at scale $y$ needs
$\eta=O(y^{-1/2+\epsilon})$.

\subsection*{Conditional logarithmic means}
For the same fixed $m$, put
\[
 \begin{gathered}
 U_m(x)=\sum_{\substack{p\le x\text{ good}\\p\equiv1\bmod m}}
        t(p)\log p(1-\sqrt{p/x}),\\
 d_m=\ord_{1/2}J_m,
 \qquad G_m=(\Z/m\Z)^\times.
 \end{gathered}
\]
Under RH for the twists in $J_m$,
\[
 \lim_{Y\to\infty}\frac1Y\int_0^Y e^{-y/2}U_m(e^y)\,dy
       =-\frac{d_m+|G_m[2]|/2}{\varphi(m)},
 \qquad G_m[2]=\{a\in G_m:a^2=1\}.
\]
Under the stronger hypothesis of full Dedekind RH for $T=EB$,
Malik and Prabhu's equidistribution theorem also gives
\[
 \lim_{X\to\infty}\frac1{\log X}\sum_{K\le X}\frac{T_m(K)}K
       =-\frac{d_m+|G_m[2]|/2}{\varphi(m)}.
\]
For $m=3$ the right-hand side is $-1/2$, since
Theorem~\ref{thm:positive} gives $d_3=0$.

\subsection*{Two quarter-power criteria for the fixed cubic function}
Let $\psi$ be either cubic character of $E/K$, with
$L_K(s,\psi)=L(s)$. For a finite set $S$ of prime ideals write
$P_S(s)=L_K^S(s,\psi)$, and restrict ideal sums to ideals coprime
to $S$. Let $\mu_K$ be the ideal M\"obius function and $N$ the
absolute ideal norm. Let $S_0$ be the primes above $3$ and $269$. Then
\[
 \sum_{\mathfrak a\perp S_0}
       \frac{\mu_K(\mathfrak a)^2\psi(\mathfrak a)}{(N\mathfrak a)^s}
        =\frac{P_{S_0}(s)}{P_{S_0}(2s)}\qquad(\Re s>1).
\]
The identity uses $L_K^{S_0}(s,\psi^2)=L_K^{S_0}(s,\psi)$.
RH for $L$ is equivalent to
\[
 \sum_{\mathfrak a\perp S_0}\mu_K(\mathfrak a)^2\psi(\mathfrak a)
       w(N\mathfrak a/X)=O_{w,\epsilon}(X^{1/4+\epsilon})
\]
for every $w\in C_c^\infty((1,2))$ and every $\epsilon>0$.

For the second criterion, set
$\eta(\mathfrak a)=\chi_{-4}(N\mathfrak a)$,
$\mathfrak l=(3+\omega)$, and let $S_1$ be the primes above
$2,3,7,269$. For each divisor $\mathfrak d$ of an odd ideal
$\mathfrak c$, take two distinguishable signs
$\eta(\mathfrak d)$ and
$\eta(\mathfrak c)\eta(\mathfrak d)$; let $W(\mathfrak c)$
count ordered pairs of opposite signs. If $\mathfrak a\perp S_1$,
\[
 \begin{gathered}
 W(\mathfrak l\mathfrak a)=8\tau_K(\mathfrak a)^2,\\
 \sum_{\mathfrak a\perp S_1}
   \frac{\psi(\mathfrak a)W(\mathfrak l\mathfrak a)}{(N\mathfrak a)^s}
       =8\frac{P_{S_1}(s)^4}{P_{S_1}(2s)}.
 \end{gathered}
\]
Here $\tau_K$ counts ideal divisors and $N\mathfrak l=7$.
RH for $L$ is also equivalent to
\[
 \sum_{\mathfrak a\perp S_1}\psi(\mathfrak a)W(\mathfrak l\mathfrak a)
       w(N\mathfrak a/X)=O_{w,\epsilon}(X^{1/4+\epsilon})
\]
for every such $w$ and every $\epsilon>0$.
In either criterion it suffices to use the following one fixed
nonnegative smooth weight. Put $L_0=\log2$, let the $U_k$ be
independent uniform variables on $[0,L_0/2^{k+1}]$, let $g$ be the
density of $L_0/4+\sum_{k\ge1}U_k$, and set $w(x)=g(\log x)$.
Its support lies in $(1,2)$, and its Mellin transform is nonzero
throughout $\Re s>0$. Section~8 proves the criteria by zero descent
and proves that this weight has the stated properties.

\subsection*{Bounds furnished by the external zero-free theorem}
For every fixed $m$ and $\epsilon>0$, Liu's named theorem implies
\[
 |P_m(y)|+|C_m(y)|\ll_{m,\epsilon}y^{B_*+\epsilon},\qquad
 T_m(K)=O_{m,\epsilon}(K^{2B_*-1+\epsilon}).
\]
For every fixed smooth compactly supported weight on $(0,\infty)$,
both ideal sums in the preceding subsection are
$O_{w,\epsilon}(X^{B_*/2+\epsilon})$.
The required square-root prime bound and quarter-power ideal bounds
remain open. The tensor, Gauss-sum, elliptic-torsion and mixed-sextic
comparisons in Sections~7--8 identify the relevant coefficients and
the costs of extracting primes; they supply no estimate at either
of these target exponents.

\section{Proof overview}
The main difficulty in real nonvanishing is to control an infinite signed theta series for a whole interval of arguments. A finite sum alone cannot do this without a tail estimate. The first 2048 partial sums of its exact integer coefficients are all at least one. Abel summation therefore makes the truncated power series at least its first monomial. A geometric tail estimate, valid simultaneously for every argument needed by the functional equation, is less than one millionth of that monomial. The theta inversion then writes the completed $L$-function at every real argument as an integral of a strictly positive function against a positive weight.

The prime criteria require a different kind of information. A smoothed prime sum has a Mellin transform proportional to $-L'/L$, so any off-line zero contributes a pole. A one-sided square-root bound makes a nonnegative Mellin transform whose first real singularity would have a negative residue, which is impossible. Prime powers and integer-square sampling have to be included with their exact costs. Unrestricted lattice cancellation removes neither the ideal M\"obius coefficients in prime extraction nor the rough composite contribution. Those are the terms for which a bound is still required.


The quarter-power criteria use the denominators of the exact ideal
Dirichlet series. If a zero off the critical line were canceled by
the numerator, the functional equation would force a sequence of
zeros converging to $2/3$. Real positivity excludes that fixed point,
and holomorphy excludes an accumulation of distinct zeros. The reverse
implication is a Mellin contour shift. This explains why the ideal
criteria have exponent $1/4$ even though prime extraction requires
exponent $1/2$.

\section{The coefficient certificate and real positivity}
This section supplies all local factors and the finite-to-infinite argument in Theorem~\ref{thm:positive}.

\subsection{Arithmetic and the completion}
The discriminant of $F$ is $-3\cdot269^2$. Here is a local derivation. The polynomial is Eisenstein at 269, giving tame total ramification of degree three and different exponent two. At three, $269$ is a cube in $\Q_3$: its negative is in $1+9\Z_3$, and cubing maps $1+3\Z_3$ bijectively onto $1+9\Z_3$ by the $3$-adic logarithm. Hence
\[
 F\otimes\Q_3\simeq\Q_3\times\Q_3(\omega),
\]
whose discriminant exponent is one. Every other prime is excluded by the polynomial discriminant $-27\cdot269^2$. The field has one complex pair, determining the negative sign. Over $K$, the cubic Hecke character has conductor $(269)$: it is nontrivially tamely ramified there, locally trivial at three by the same cube calculation, and unramified elsewhere. Its conductor norm is $269^2$.

The local coefficient polynomials are
\[
 \begin{array}{c|c}
 p&\displaystyle\sum_{k\ge0}A(p^k)T^k\\\hline
 3&(1-T)^{-1}\\
 269&1\\
 p\equiv2\pmod3,\ p\ne269&(1-T^2)^{-1}\\
 p\equiv1\pmod3&(1-t(p)T+T^2)^{-1}.
 \end{array}
 \tag{3.1}
\]
Thus the last row is $(1-T)^{-2}$ if $t(p)=2$ and $(1+T+T^2)^{-1}$ if $t(p)=-1$. These formulas show multiplicativity, integrality, and
\[
 |A(n)|\le d(n)\le n.
 \tag{3.2}
\]
Here $d(n)$ is the number of positive divisors of $n$. For example
\[
 A(1),\ldots,A(13)=1,0,1,1,0,0,-1,0,1,0,0,1,-1.
\]

The signatures of $F$ and $\Q$ are $(1,1)$ and $(1,0)$. Dividing their completed Dedekind zeta functional equations cancels the real gamma factor and gives (1.1), including root number $+1$. Induction from the nontrivial finite-order Hecke character supplies entireness. Mellin inversion of the gamma factor, followed by shifting a contour across an entire completed function, gives
\[
 f(t)=t^{-1}f(1/t).
 \tag{3.3}
\]
One can justify the contour shift first in a fixed vertical strip: the gamma factor has exponential decay and the finite-order Hecke function has polynomial growth there. There are no poles to cross. Consequently
\[
 \frac{\Lambda(s)}2
  =\int_1^\infty f(t)\bigl(t^{s-1}+t^{-s}\bigr)\,dt
 \qquad(s\in\C).
 \tag{3.4}
\]
Both integrals converge locally uniformly in $s$, because $f(t)$ decays exponentially as $t\to\infty$.

\subsection{Independent integer computations}
There is also a direct finite-field construction of the coefficients. The prime 269 is inert in $\Z[\omega]$, so
\[
 \Z[\omega]/(269)=\mathbb F_{269^2}.
\]
For $v\ne0$ let $\xi(v)$ be the complex cube root of unity corresponding to $v^{24120}$; extend it by zero at zero. It kills all six Eisenstein units, is trivial on rational units, and is inverted by conjugation. The Eisenstein ring has class number one, so its ray class group modulo $(269)$ is
\[
 \mathbb F_{269^2}^{\times}/\mu_6,
\]
a cyclic group of order $12060$ (Milne, 2020, Theorem V.1.7, p.~150). It has precisely two nontrivial cubic characters, which are conjugate. Since $E/K$ has degree three and conductor $(269)$, the character in (1.1) is $\xi$ or $\bar\xi$, by the reciprocity and existence theorems (Milne, 2020, V.3.5--3.7, p.~158). Conjugation permutes ideals of each norm, so their rational Dirichlet coefficients coincide. Therefore
\[
 A(n)=\frac1{12}\sum_{a^2-ab+b^2=n}
       \bigl(\xi(a+b\omega)+\overline{\xi(a+b\omega)}\bigr).
 \tag{3.5}
\]
The summand is $2$ when the residue is a nonzero cube, $-1$ when it is a noncube, and zero at the zero residue. This computes $A(n)$ using only arithmetic modulo 269 and enumeration of lattice points. The factor 12 is six associates times the sum of a character and its conjugate. At the prime three, $(1-\omega)^2=-3\omega$ implies $\xi(1-\omega)=1$, agreeing with (3.1).

The two methods agree at every $1\le n\le8192$. For $N=2048$, their cumulative sums $S_k=\sum_{n\le k}A(n)$ give
\[
 \min_{1\le k\le N}S_k=1,\qquad
 \{k:S_k=1\}=\{1,2\},\qquad S_N=47.
 \tag{3.6}
\]
The \emph{Cubic theta certificate} program in the \href{\workbenchurl}{workbench} contains both algorithms, checks every equality and minimum, and evaluates the exact integer inequality used below. It needs no numerical library for (3.6) or the positivity proof. The optional central-value calculation uses Arb interval arithmetic.

\subsection{A uniform rational tail}
Put $x=\exp(-2\pi t/\sqrt Q)$. For $t\ge1$ we have $x<75/76$: indeed $\sqrt Q<466$ and $\pi>25/8$ imply $2\pi/\sqrt Q>1/75$, while $e^u>1+u$ gives $e^{-1/75}<75/76$. Abel summation and (3.6) give
\[
 \sum_{n=1}^N A(n)x^n
  =(1-x)\sum_{k=1}^{N-1}S_kx^k+S_Nx^N\ge x.
 \tag{3.7}
\]
Using (3.2), the absolute tail divided by $x$ is at most
\[
 \sum_{n>N}n x^{n-1}
 \le 76(N+76)\left(\frac{75}{76}\right)^N
 <10^{-6}.
 \tag{3.8}
\]
The final inequality is the integer comparison
\[
 10^6\cdot76\cdot2124\cdot75^{2048}<76^{2048}.
\]
Equations (3.7)--(3.8) prove (1.2). For real $s$, the weight in (3.4) is positive. Hence $\Lambda(s)>0$. The factor relating it to $L(s)$ in (1.1) is positive for $s>0$, proving real nonvanishing.

At $s=1/2$, (3.4) gives the particularly convenient formula
\[
 L(1/2)=2\sum_{n\ge1}\frac{A(n)}{\sqrt n}
                  \operatorname{erfc}\sqrt{\frac{2\pi n}{\sqrt Q}}.
 \tag{3.9}
\]
Writing $a=2\pi/\sqrt Q$, the tail after $M$ terms has absolute value at most
\[
 \frac{2e^{-a(M+1)}}{\sqrt{\pi a}(1-e^{-a})}.
 \tag{3.10}
\]
This follows from $\operatorname{erfc}(u)\le e^{-u^2}/(\sqrt\pi u)$ and $|A(n)|\le n$. For $M=8192$, the bound is less than $7.531\cdot10^{-46}$. Computing the finite sum with 240-bit ball arithmetic and enlarging by this tail gives (1.3).

\section{The exact signed-square criterion}
This section derives the arithmetic inequality at the square-root threshold, including its real residue and sampling error.
Fix $3\mid m$ and let $\chi$ range over characters modulo $m$, interpreted as complete primitive Artin factors. Put
\[
 J_m(s)=\prod_{\chi\bmod m}L(s,\rho\otimes\chi),\qquad
 R_m(x)=\sum_{\substack{p\le x\text{ good}\\p\equiv1\bmod m}}t(p)\log p,
 \quad T_m(K)=\frac1{K^2}\sum_{k\le K}R_m(k^2).
 \tag{4.1}
\]
Every factor is an entire nontrivial Hecke $L$-function over $K$. In fact a norm pullback is fixed by conjugation, whereas the cubic character is inverted, so the induced twist cannot become trivial. If $B=\Q(\zeta_m)$ and $T=EB$, Artin formalism gives $J_m=\zeta_T/\zeta_B$. This identity does not require any Riemann hypothesis.

\begin{proof}[Proof of Theorem~\ref{thm:squares}]
Write $-J_m'/J_m=\sum b_m(n)n^{-s}$ and
\[
 A_m(x)=\sum_{n\le x}b_m(n)(1-\sqrt{n/x}).
\]
At good primes, $b_m(p)=\varphi(m)t(p)\log p$ if $p\equiv1\pmod m$ and is zero otherwise. For all prime powers, $|b_m(p^j)|\le2\varphi(m)\log p$. Chebyshev's elementary estimate gives an $O_m(\sqrt x)$ contribution from squares, an $O_m(x^{1/3}\log(2x))$ contribution from higher powers, and an $O_m(\log(2x))$ contribution from the excluded primes. Thus all omitted terms cost $O_m(\sqrt x)$.

A prime $p\le K^2$ occurs in $K-\lfloor\sqrt p\rfloor$ samples. The rounding cost is bounded by
\[
 K^{-1}\sum_{p\le K^2}2\varphi(m)\log p=O_m(K).
\]
Hence
\[
 A_m(K^2)=\varphi(m)K T_m(K)+O_m(K).
 \tag{4.2}
\]
The same coefficient bounds imply $\sum_{n\le x}|b_m(n)|=O_m(x)$. If $K\le\sqrt x<K+1$, changing the cutoff weight changes $A_m$ by $O_m(K)$; the weight changes by $O(1/K)$ and the coefficient mass is $O_m(K^2)$. Therefore the lower bound for $T_m$ is equivalent to $A_m(x)\ge-C_m\sqrt x$ for all $x\ge1$.

For $\Re s>1$, termwise integration gives
\[
 \int_1^\infty A_m(x)x^{-s-1}\,dx
   =-\frac{J_m'(s)}{J_m(s)}W(s),\qquad
 W(s)=\frac1{s(2s+1)}.
 \tag{4.3}
\]
Apply Landau's positivity principle to $A_m(x)+C_m\sqrt x$. Its transform is (4.3) plus $C_m/(s-1/2)$. If its abscissa of convergence were greater than $1/2$, that real point would have to be a singularity. A real zero $\beta>1/2$ of order $h$ would give residue
\[
                    -\frac{h}{\beta(2\beta+1)}<0.
 \tag{4.4}
\]
A nonnegative Mellin integral cannot tend to minus infinity when approached from the right. At one the logarithmic derivative is regular, because $J_m(1)\ne0$. These facts force the abscissa to be at most $1/2$, and holomorphy of the nonnegative integral excludes every zero with real part greater than $1/2$. The functional equations, with the conjugate twists also in the family, reflect all left-hand zeros to the right.

Conversely, assume RH for the twists. Mellin inversion of (4.3), shifting from $\Re s=2$ to $\Re s=1/4$, gives
\[
 A_m(x)=-\sum_\rho \ord_\rho(J_m)W(\rho)x^\rho+O_m(x^{1/4}).
 \tag{4.5}
\]
The sum converges absolutely since $W(1/2+it)=O((1+|t|)^{-2})$ and the number of zeros in a unit height interval is $O_m(\log(2+|t|))$. Horizontal contours can be chosen at distance at least a constant times $\log^{-2}(2+|t|)$ from zero ordinates; the usual local partial fraction formula bounds $J_m'/J_m$ there by $O_m(\log^3(2+|t|))$. Their integrals vanish, and on $\Re s=1/4$ the vertical integral is $O_m(x^{1/4})$. These are the standard fixed-Hecke explicit-formula estimates (Iwaniec \& Kowalski, 2004, Chapter 5). Thus $A_m(x)=O_m(\sqrt x)$ and (4.2) proves the asserted bound.
\end{proof}

At $m=3$, the two twists coincide, so $J_3=L^2$. Theorem~\ref{thm:positive} excludes real zeros of this product without assuming Theorem~\ref{thm:squares}. At $m=12$,
\[
 J_{12}=L(s,\rho)^2L(s,\rho\otimes\chi_{-4})^2.
\]
The Gaussian twist is tested individually by replacing $R_m$ by $R_{12,1}-R_{12,7}$ in the same proof.

\subsection{Continuous and sampled biases}
The normalization of the average has a computable bias. Put
\[
 U_m(x)=\sum_{\substack{p\le x\text{ good}\\p\equiv1\bmod m}}
     t(p)\log p(1-\sqrt{p/x}),\quad
 d_m=\ord_{1/2}J_m,
 \quad G_m=(\Z/m\Z)^\times.
\]
Under RH for the twists, the logarithmic mean is
\[
 \lim_{Y\to\infty}\frac1Y\int_0^Y e^{-y/2}U_m(e^y)\,dy
  =-\frac{d_m+|G_m[2]|/2}{\varphi(m)}.
 \tag{4.6}
\]
Here $G_m[2]=\{a:a^2=1\}$. To prove the prime-square term, average
$\tr\rho(g^2)\tr\operatorname{reg}_{G_m}(a^2)$ over the compatible pairs in $S_3\times G_m$. Only $a^2=1$ contribute. Exactly half of these elements act nontrivially on $K$, by pairing with $-1$. Above them $g$ is a transposition and $\tr\rho(g^2)=2$; above the others the three traces are $2,-1,-1$ and sum to zero. The average coefficient of $\log p$ at $p^2$ is $|G_m[2]|$. Chebotarev and partial summation therefore give $\tfrac12|G_m[2]|\sqrt x+o_m(\sqrt x)$ for their smoothed contribution. Higher prime powers are $o_m(\sqrt x)$. In (4.5), nonzero frequencies have zero logarithmic mean, while $W(1/2)=1$, proving (4.6).

The exact rounding identity is
\[
 T_m(K)-\frac{U_m(K^2)}K
 =\frac1{K^2}\sum_{\substack{p\le K^2\text{ good}\\p\equiv1\bmod m}}
                  t(p)\{\sqrt p\}\log p.
 \tag{4.7}
\]
Under full Dedekind RH for $T$, Malik and Prabhu (2022, Theorem 1.1) with $\alpha=1$, $\theta=1/2$, $\Delta=1/2$ and auxiliary parameter $\omega=\log x$ gives joint equidistribution of $\{\sqrt p\}$ in each fixed Frobenius class. Its errors are $o(x/\log x)$; partial summation gives $o(x)$ after logarithmic weighting. Approximation of $u\in[0,1]$ by step functions then makes the right side of (4.7) tend to one half the mean trace in the relevant fibre, which is zero. Absolute convergence of (4.5), and the boundedness of $\sum_{k\le X}k^{-1+it}$ for fixed nonzero $t$, yield
\[
 \lim_{X\to\infty}\frac1{\log X}\sum_{K\le X}\frac{T_m(K)}K
       =-\frac{d_m+|G_m[2]|/2}{\varphi(m)}.
 \tag{4.8}
\]
For $m=3$, Theorem~\ref{thm:positive} gives $d_3=0$, so this conditional sampled logarithmic mean is exactly $-1/2$.

\section{Prime windows, their centering, and the identified period}
This section states the precise quantitative target and compares it with a single logarithmic theta moment.
Fix $\delta>0$ and define
\[
 P_m(y)=R_m(2y)-R_m(y).
\]
Alp\"oge's general Landau-window theorem \cite{AlpogeWindow} implies that the twist RH in Theorem~\ref{thm:squares} is equivalent to the following assertion: for every $0<\epsilon<1/2$, there is $C_{m,\epsilon,\delta}$ and arbitrarily large $Y$ such that
\[
 P_m(y)\ge-C_{m,\epsilon,\delta}y^{1/2+\epsilon}
       \quad\hbox{for every }Y\le y\le Y^{1+\delta}.
 \tag{5.1}
\]
All parameters are fixed before the arbitrarily large windows are chosen; no uniformity in $m$ is asserted. For the forward reduction set $g(u)=-\sum_{e^u<n\le2e^u}b_m(n)$. Its Laplace transform is
\[
 \int_0^\infty g(u)e^{-su}\,du
       =\frac{2^s-1}{s}\frac{J_m'(s)}{J_m(s)}.
 \tag{5.2}
\]
The multiplier is nonzero in $\Re s>0$. At a positive real zero $\beta$ of order $h$, its residue is $h(2^\beta-1)/\beta>0$. The prime-power error is $O_m(\sqrt y\log^2(2y))$, which is smaller than the right side of (5.1) for each fixed positive $\epsilon$. The majorant has weighted integral $O_{m,\epsilon,\delta}(\log Y)$ on $[\log Y,(1+\delta)\log Y]$, satisfying the analytic theorem. Conversely the fixed-Hecke explicit formula under RH gives $P_m(y)=O_m(\sqrt y\log^2(2y))$ for every $y$.

At the exact square-root threshold, an equivalent averaged test uses the negative part $v_-=\max(0,-v)$:
\[
 \sum_{K\le k\le K^{1+\delta}}\frac{(T_m(k))_-}{k}\le K^\epsilon
 \tag{5.3}
\]
for every $\epsilon>0$ on arbitrarily late integer windows. Apply the same analytic theorem to $-A_m(e^u)$ and use (4.2) with $u=2\log k$. The interpolation error contributes $O_{m,\delta}(\log K)$, and choosing the exponent in (5.3) smaller than twice the target Laplace loss gives its weighted majorant. This specifies why (5.3) uses an arbitrarily small power even though $T_m$ itself is normalized at exponent exactly $1/2$.

The ramp form has an additional centering issue. Set $D_m=J_m'/J_m$, let $b_{m,1}$ be the part of $b_m$ supported on first powers of good primes, and put $h_m=b_m-b_{m,1}$. Define
\[
 C_m=D_m(1)+\sum_{n\ge1}\frac{h_m(n)}n,
 \quad
 V_m(y)=\sum_{p\le y}b_{m,1}(p)(y/p-1)+yC_m.
 \tag{5.4}
\]
The correction converges absolutely. The full ramp and its transform are
\[
 \begin{split}
 G_m(y)&=\sum_{n\le y}b_m(n)(y/n-1)+yD_m(1),\\
 \int_1^\infty G_m(y)y^{-s-1}\,dy
     &=\frac{sD_m(1)-D_m(s)}{s(s-1)}.
 \end{split}
 \tag{5.5}
\]
At a zero $z$ the residue is $\ord_z(J_m)/(z(1-z))$, positive for real $0<z<1$. Direct subtraction, including the omitted tail, gives
\[
 G_m(y)-V_m(y)
    =-\sum_{n\le y}h_m(n)-y\sum_{n>y}\frac{h_m(n)}n
    =O_m(\sqrt y\log^2(2y)).
 \tag{5.6}
\]
Thus an upper bound $V_m(y)\le C y^{1/2+\epsilon}$ on the same windows is the corresponding ramp criterion. Replacing $C_m$ by an approximation with error $\eta$ changes the ramp by $\eta y$; to test a scale $y$, its error must be $O(y^{-1/2+\epsilon})$. A fixed decimal value of the centering constant cannot certify unbounded windows.

For $m=3$, $D_3(1)=2L'(1)/L(1)$. The concrete ray-theta period identifies this logarithmic derivative exactly. With the unscaled lattice theta
\[
 \Theta_\xi(v)=\sum_{\alpha\in\Z[\omega]}\xi(\alpha)e^{-\pi v|\alpha|^2},
 \quad I(s)=\int_0^\infty\Theta_\xi(v)v^{s-1}\,dv,
\]
six associates per ideal give $I(s)=6\pi^{-s}\Gamma(s)L(s)$ and
\[
 \frac{L'(1)}{L(1)}=\log\pi+\gamma+
       \frac{\int_0^\infty\Theta_\xi(v)\log v\,dv}
            {\int_0^\infty\Theta_\xi(v)\,dv}.
 \tag{5.7}
\]
The theta transformation proves absolute convergence at both ends. The denominator is positive by the residue formula for $\zeta_F/\zeta$; Theorem~\ref{thm:positive} also makes the appropriately scaled kernel positive. Formula (5.7) identifies a scalar. Inequalities (5.1), (5.3), and (5.4) ask for a family of prime or Euler-deletion estimates at unbounded scales. No inequality of that strength follows from (5.7) alone.

\section{What squarefree composites and bilinear splitting leave}
This section quantifies the range of factors for which an additional arithmetic estimate is necessary.
Let $z=(\log y)^2$, and let $C_m(y)$ be the sum of
\[
 \left(\prod_{p\mid n}t(p)\right)\log n
\]
over squarefree $y<n\le2y$ with at least two prime factors, $n\equiv1\pmod m$, and every prime factor good and greater than $z$. Then, unconditionally for every fixed $m$ and $\epsilon>0$,
\[
 P_m(y)+C_m(y)=O_{m,\epsilon}(y^{1/2+\epsilon}).
 \tag{6.1}
\]
Here is the coefficient proof, to make the quantifiers and the unhandled part explicit. Periodic cancellation of $\xi$ on the Eisenstein lattice gives
\[
 \sum_{\substack{n\le X\\ n\equiv b\bmod m}}A(n)=O_m(\sqrt X)
 \tag{6.2}
\]
for every reduced $b$. To see this, choose modulo $\operatorname{lcm}(269,m)$ a unit $v$ whose cubic character is nontrivial, and multiply each norm fibre by $v/\bar v$. This preserves its norm and multiplies its character sum by $\xi(v)^2\ne1$, so a full period contributes zero. Only period cells meeting the boundary of the norm disk remain, and their count is $O_m(\sqrt X)$. The ray-class identification in Section~3.2 supplies the ideal character used here.

We will also need the uniform, slightly weaker bound $O_\epsilon(X^{1/2+\epsilon})$ when $m$ varies. Fix a rational integer $C$ for the ray character's period and tile by $Cm\Z[\omega]$. For $m\le\sqrt X$, there are $O_C(\sqrt X/m+1)$ boundary tiles. A tile contains at most $C^2$ representatives of each residue modulo $m$ with norm $b$. The number of those residues is $O(m\log(2m))$: a nonempty unit norm fibre modulo $p^k$ has size $\varphi(p^k)$ for split $p$, size $p^{k-1}(p+1)$ for inert $p$, and at most $2\cdot3^k$ at three. The last assertion follows since the norm image contains the squares of rational units, of index two. Multiplication of these bounds and $\prod_{p\mid m}(1+1/p)\le\sum_{d\mid m}1/d\ll\log(2m)$ gives the claim. The boundary contribution is therefore $O_C(\sqrt X\log(2X))$. For $m>\sqrt X$ there are at most $1+X/m$ possible integer norms, each with at most $d(n)\ll_\epsilon n^\epsilon$ ideals. This proves the uniform bound. The argument also applies to the mixed ray character below: at the factor above 269 choose $v$ with nontrivial cubic character and at the other conductor factors choose $v=1$.

In the Euler product $F_z(s)=\prod_{p>z,\,p\text{ good}}(1+t(p)p^{-s})$, write $F_z=L H_z$. The local factors of $H_z$ at large good primes are
\[
 1-3T^2+2T^3,\quad1-T^3,\quad1-T^2
 \tag{6.3}
\]
for identity, three-cycle, and transposition respectively. At small primes one uses the reciprocal local polynomial in (3.1). For $1/2<\sigma<1$ the absolute coefficient norm is at most
\[
 \sum_d|h_z(d)|d^{-\sigma}
 \ll_{m,\sigma}\exp\left(C_\sigma z^{1-\sigma}/\log(2z)\right).
 \tag{6.4}
\]
The large-prime product converges because its linear terms vanish; the small-prime estimate is Chebyshev and partial summation. Convolve (6.2) with $h_z$, use $(X/d)^{1/2}\le(X/d)^\sigma$, and insert $\log n$ by partial summation. For $\sigma=1/2+u$ and any fixed $u>0$, the exponential in (6.4) is $y^{o(1)}$. Taking $u<\epsilon$ proves (6.1). This argument also allows any $z=(\log y)^{2+o(1)}$.

Consequently the missing prime lower bound is exactly the upper bound
\[
 C_m(y)\le C_{m,\epsilon,\delta}y^{1/2+\epsilon}
 \tag{6.5}
\]
throughout the arbitrarily late windows of (5.1). The sign reverses because of (6.1).

The number of distinct factors is not bounded independently of $y$:
\[
 \omega(n)\le\frac{\log(2y)}{\log z}
       =\left(\frac12+o(1)\right)\frac{\log y}{\log\log y}.
 \tag{6.6}
\]
For every fixed $k$, sufficiently large scales have $k$ distinct split primes in $(y^{1/k},2^{1/k}y^{1/k})$, by the prime number theorem in the progression; their product lies in $(y,2y]$ and each exceeds $z$. Thus a fixed-semiprime calculation omits infinitely many eventual factor types.

One can expose a balanced part without losing any terms. In a squarefree product whose every prime factor is at most $y^{1/6}$, multiply its ordered prime factors until the first product $u$ exceeds $y^{1/3}$. Then
\[
 y^{1/3}<u\le y^{1/2},\qquad y^{1/2}<v=n/u<2y^{2/3}.
 \tag{6.7}
\]
The split is unique, has $(u,v)=1$, and requires the ordering condition that every prime factor of $u$ precedes those of $v$. A bilinear estimate must retain that condition and the roughness and character weights. The complement consists of products having a prime factor greater than $y^{1/6}$, including products $ap$ with $z<a\le y^{1/3}$ and $p\asymp y/a$. They have not been estimated by (6.2). Even before this split, absolute values give only $y^{1+o(1)}\log y$: by (6.6), each product weight is at most $2^{\omega(n)}=y^{o(1)}$, and there are at most $2y$ integers. The gap from (6.5) is a half power of the scale.

There is also a structural restriction on a prospective bilinear theorem. On a rectangular pair of sets of ideals prime to the conductor, multiplicativity gives
\[
 \sum_{\mathfrak a\in\mathcal A,\,\mathfrak b\in\mathcal B}
    \lambda_{\mathfrak a}\mu_{\mathfrak b}\xi(\mathfrak a\mathfrak b)
 =\left(\sum_{\mathfrak a\in\mathcal A}\lambda_{\mathfrak a}\xi(\mathfrak a)\right)
  \left(\sum_{\mathfrak b\in\mathcal B}\mu_{\mathfrak b}\xi(\mathfrak b)\right).
\]
Choosing the bounded coefficients to be the conjugate character makes this equal $|\mathcal A||\mathcal B|$. Thus an estimate uniform over arbitrary bounded coefficient arrays cannot save a positive power. A successful bilinear argument must use the particular prime-extraction coefficients or introduce additional oscillation with a proved identity. Merikoski's cubic spin over primes has its own reciprocity weight; his result does not replace this finite-conductor character by that weight.

The obstruction can also be written with the ideal M\"obius function. If $A^{-1}$ denotes the Dirichlet convolution inverse, then
\[
 \Lambda_\rho=(A\log)*A^{-1},\qquad
 A^{-1}(n)=\sum_{N\mathfrak d=n}\mu_K(\mathfrak d)\xi(\mathfrak d).
 \tag{6.8}
\]
Using (6.2) on the part $N\mathfrak d\le D$ gives
\[
 O_m(\sqrt{yD}\log(2y)).
 \tag{6.9}
\]
Indeed the inner sum is $O_m(\sqrt{y/N\mathfrak d}\log(2y))$ and the number of ideals of norm at most $D$ is $O(D)$. A polylogarithmic $D$ is admissible for the target; the complementary inverse sum with $D<N\mathfrak d\le2y$ is the precise uncontrolled term. Periodic cancellation without its M\"obius weight does not bound it.

\section{Positivity and averaging do not estimate the centered excess}
This section identifies what ordinary positivity, parameter averaging, and the available metaplectic estimates control, together with the estimates still needed.

\subsection{Rankin--Selberg positivity}
The $S_3$ character table gives
\[
 \rho\otimes\bar\rho\simeq\mathbf1\oplus\eta\oplus\rho,
 \qquad L(s,\rho\otimes\bar\rho)=\zeta(s)L(s,\eta)L(s,\rho),
 \tag{7.1}
\]
where $\eta=\chi_{-3}$. This is the Rankin--Selberg square of the dihedral weight-one form. At every unramified prime power its logarithmic-derivative coefficient is $\log p\,|\tr\rho(\Frob_p^k)|^2\ge0$. In the prime window at $m=3$, put
\[
 Q_2(y)=\sum_{\substack{y<p\le2y\\p\text{ good}}}t(p)^2\log p,
 \quad W_1(y)=\sum_{\substack{y<p\le2y\\p\equiv1\bmod3,\ p\text{ good}}}\log p.
\]
Because $t^2=t+2$ on the two split classes,
\[
                         P_3(y)=Q_2(y)-2W_1(y).
 \tag{7.2}
\]
The positivity $Q_2\ge0$ only gives $P_3\ge-2W_1$; even $Q_2\ge W_1$ only recovers the elementary bound $t\ge-1$. Both are of order $y$. The desired result is precisely $Q_2\ge2W_1-O(y^{1/2+\epsilon})$, an estimate of the centered positive sum. Equation (7.1) retains the same unknown factor $L(s,\rho)$, with multiplicity one. Thus ordinary Rankin--Selberg positivity supplies no exponent improvement here.

More generally, for every integer $n\ge1$,
\[
 \rho^{\otimes n}=a_n(\mathbf1\oplus\eta)\oplus b_n\rho,
 \quad a_n=\frac{2^n+2(-1)^n}{6},\quad
 b_n=\frac{2^n-(-1)^n}{3}.
 \tag{7.3}
\]
This follows by comparing the characters $2^n,0,(-1)^n$ on the three classes. The centered prime sum on $1\pmod m$ is exactly $b_nR_m$, so taking ordinary tensor powers only multiplies the same unresolved statistic.

\subsection{Positive Mellin kernels}
Theorem~\ref{thm:positive} also does not prove RH by an abstract positivity principle. Here is an explicit obstruction to that inference. Let $k$ be any nonzero nonnegative smooth even function of compact support, let $0<c<1$, and set
\[
 h(u)=k(u)+\frac c2\bigl(k(u-a)+k(u+a)\bigr).
\]
Its bilateral Laplace transform is
\[
 H(z)=K(z)(1+c\cosh(az)),\qquad K(z)=\int_\R k(u)e^{zu}\,du.
\]
The kernel is nonnegative and even, $H$ is positive on the real axis, and $H(z)=H(-z)$. Yet $H$ has zeros at
\[
 z=\frac{\pm\operatorname{arcosh}(1/c)+(2j+1)\pi i}{a}.
 \tag{7.4}
\]
With $a=4$ and $c=1/2$, the real parts of $1/2+z$ lie strictly between zero and one and are not $1/2$. Multiplication by $K$ cannot remove these zeros. This does not construct an Artin $L$-function with off-line zeros; it proves that real positivity and reflection of a Mellin kernel, without an additional arithmetic property, cannot establish the desired complex-zero statement.

\subsection{Averaging the radicand}
For $p\equiv1\pmod3$ and variable radicand $a$, Fourier inversion gives
\[
 t_a(p)=\frac1p\sum_{h=1}^{p-1}e^{-2\pi iha/p}
                    \sum_{x\bmod p}e^{2\pi ihx^3/p}.
\]
There are $(p-1)/3$ nonzero residues of weight two and $2(p-1)/3$ of weight minus one, so their mean is zero and their squared mass is $2(p-1)$. Distinct fractions $h/p$ with $y<p\le2y$ are separated by at least $1/(4y^2)$. The dual additive large sieve therefore bounds the squared mean of $T_a=\sum t_a(p)\log p$ over $1\le a\le H$ by
\[
 \sum_{a\le H}|T_a|^2\ll(H+4y^2)y\log(2y).
 \tag{7.5}
\]
At $H=y^2$ this gives square-root size on average. The copies $a=269b^3\le H$ have the same trace at each prime once $p>b$. There are only $\asymp y^{2/3}$ such copies, so (7.5) gives for the prescribed radicand only
\[
 |T_{269}|\ll y^{7/6}\sqrt{\log(2y)},
 \tag{7.6}
\]
weaker than the elementary estimate of order $y$. This computes, rather than suppresses, the cost of passing from the family mean to 269.

Recent cubic family moment and zero-density results may be valuable inputs to a future argument. Their family quantifiers do not themselves control a prescribed conductor and all heights. For example, a bound for the number of exceptional parameters up to $P$ by $P^\alpha$, $\alpha>0$, permits a fixed parameter such as 269 to remain exceptional at every $P$. The recent rStation papers on mixed cubic moments and on smooth amplifiers explicitly leave a complementary signed inverse term unbounded. We use none of their conclusions as an individual RH theorem.

\subsection{Known zero-free boundaries and the remaining gap}
OpenAI (2026, Theorem 1.1, p.~4) applies to every fixed twist in $J_m$: each is induced from a finite-order Hecke character of $K$. This gives the following quantitative background bound, which is weaker than the required square-root estimate:
\[
 |P_m(y)|+|C_m(y)|\ll_{m,\epsilon}y^{7/8+\epsilon}.
 \tag{7.7}
\]
Here is the implication from that external zero-free theorem. At a half-integer $x$ and a height $T$ between $x^2$ and $2x^2$ avoiding zero ordinates, the fixed-Hecke truncated explicit formula gives
\[
 \sum_{n\le x}b_m(n)
   =-\sum_{|\Im\rho|\le T}\ord_\rho(J_m)\frac{x^\rho}{\rho}
      +O_m(\log^2(2x)).
\]
The usual error $O_m(x\log^2(2xT)/T+\log^2(2x))$ has been simplified with this choice of $T$. The zero count and $\Re\rho\le7/8$ bound the sum by $O_m(x^{7/8}\log^2(2x))$. Passing to a neighbouring half-integer costs $O_m(\log(2x))$. Removing the higher prime powers costs $O_m(\sqrt x\log^2(2x))$, and taking a dyadic difference proves the prime part of (7.7). Equation (6.1) proves the composite part. Every constant is allowed to depend on the fixed modulus $m$. The same proof with Liu's main boundary gives $|P_m(y)|+|C_m(y)|\ll_{m,\epsilon}y^{B_*+\epsilon}$. The remaining exponent gap is $B_*-1/2=0.3733458561067\ldots$. In the square statistic this gives only $T_m(K)=O_{m,\epsilon}(K^{2B_*-1+\epsilon})$, rather than a bounded negative part.

\subsection{The exact Patterson coefficient correlation}
The cubic metaplectic theta series does encode the fixed radicand through two Fourier coefficients. We give that identity before comparing the available estimates. Put $r=-269$, which is primary in the convention $c\equiv1\pmod3$, and put $\lambda=\sqrt{-3}$. For a primary squarefree $c$, write
\[
 G(a,c)=\frac1{\sqrt{Nc}}\sum_{v\bmod c}\left(\frac vc\right)_3
       \exp\bigl(2\pi i\tr_{K/\Q}(av/c)\bigr),\qquad G(c)=G(1,c).
\]
If $(a,c)=1$, changing variables gives $G(a,c)=\chi_c(a)^{-1}G(c)$, and $|G(c)|=1$. Chinese remainders and cubic reciprocity give
\[
 G(rc)=G(r)G(c)\overline{\chi_c(r)},\qquad
 G(c)\overline{G(269,c)}=\chi_c(269)
       =G(r)G(c)\overline{G(rc)}.
 \tag{7.8}
\]
The identity holds with value zero on both sides if a nonsquarefree $c$ is allowed, provided the character term is then multiplied by $\mu_K(c)^2$. Here and below $c$ is prime to $3r$.

In fact $G(r)=1$. The cubic character on $\mathbb F_{269^2}^{\times}$ is trivial on $\mathbb F_{269}^{\times}$. If $d$ spans the trace-zero line, then $d^2\in\mathbb F_{269}^{\times}$, so $\chi(d)^2=1$ and $\chi(d)=1$. Thus the sum of $\chi$ over the trace-zero fibre is 268. The other 268 trace fibres have a common sum, by rational scaling; the total character sum is zero, so their common sum is $-1$. Summing against the nontrivial trace additive character gives the unnormalized value 269, hence $G(r)=1$. The accompanying exact program checks every trace fibre and separately expands (7.8) in $\Z[\omega][\Z/p\Z]$ at fourteen split rational primes; the proof above applies at every prime.

Let $\psi(c)=\chi_c(269)=\chi_r(c)$. This is one of the two cubic ray characters already identified in Section~3.2, and both have the same $L$-function. For the prime correlation define
\[
 \mathcal C_r(y)=\sum_{\substack{y<N\pi\le2y\\\pi\ {\rm primary\ prime},\ (\pi,3r)=1}}
              G(\pi)\overline{G(r\pi)}\log N\pi.
\]
Conjugate prime ideals over a split rational prime contribute its real trace $t(p)\log p$. An inert prime ideal has norm $p^2$ and character one on 269. Consequently the exact relation is
\[
 \mathcal C_r(y)=P_3(y)+
  2\sum_{\substack{\sqrt y<p\le\sqrt{2y}\\p\equiv2\bmod3,\ p\ne269}}\log p
       =P_3(y)+O(\sqrt y).
 \tag{7.9}
\]
Thus a bound $|\mathcal C_r(y)|\ll_\epsilon y^{1/2+\epsilon}$ would prove the desired RH. The one-sided power-window version of Section~5 is already sufficient. This identifies a particular two-coefficient estimate, with its exact correction, rather than merely comparing two different automorphic functions.

There is an unconditional estimate at this size after the prime restriction is removed:
\[
 \sum_{\substack{Nc\le X\\(c,3r)=1}}G(c)\overline{G(rc)}
  =\sum_{\substack{Nc\le X\\(c,3r)=1}}\mu_K(c)^2\psi(c)
  \ll\sqrt X\log(2X).
 \tag{7.10}
\]
Indeed periodicity gives $\sum_{Nc\le Y,(c,3r)=1}\psi(c)=O(\sqrt Y)$. Insert $\mu_K(c)^2=\sum_{d^2\mid c}\mu_K(d)$ and sum $\sqrt X/Nd$ over $Nd\le\sqrt X$. The number of ideals of norm at most $T$ is $O(T)$, so the resulting harmonic sum is $O(\log(2X))$. This proof uses the fixed character and gives a useful full-sum bound, but the prime restriction remains substantive.

There is also a sharper target for these unrestricted squarefree correlations. Replacing their square-root bound by $O_{W,\epsilon}(X^{1/4+\epsilon})$ for every smooth dyadic $W$ is equivalent to RH. The Dirichlet series is $P(s)/P(2s)$, and the zero-descent argument in Section~8.4 applies with numerator exponent one in place of four. A zero $\rho$ to the right of the half line would force zeros at $\rho/2$ and then $1-\rho/2$, eventually accumulating at $2/3$. Real positivity excludes the fixed point. The one fixed weight constructed there also suffices. Thus the elementary full-sum argument (7.10) has a quarter-power gap for this alternative, while (7.13) records the prime-extraction gap.

The Fourier normalization also makes the cube completion explicit. Mixed shifted coefficient sums in the Rankin--Selberg setting are discussed by de Faveri, Dunn, and Hoffstein (2026, Section 2.2, p.~7). The following is a fixed-$r$ specialization of Patterson's known cubic coefficients; no new general correlation theorem is asserted. Let $D(x)=\overline{\tau(\lambda^{-3}x)}$, where $\tau$ is Patterson's coefficient in Dunn--Radziwi\l\l\ (2024, arXiv v3, (5.7), p.~17). On primary indices,
\[
 D(nb^3)=3^{5/2}|b|G(n)\quad(n\ {\rm squarefree}),
\]
with no requirement that $(n,b)=1$, and all other primary indices have zero coefficient. Put $S=\{(\lambda),(r)\}$, $P(s)=L_K^{(S)}(s,\psi)$, and $Z_K^{(S)}(s)=\zeta_K^{(S)}(s)$. Since $G(r)=1$, the exact coefficient correlation has the series
\[
 \frac1{3^5}\sum_{\substack{x\ {\rm primary}\\(x,S)=1}}
       \frac{D(x)\overline{D(rx)}}{Nx^s}
   =Z_K^{(S)}(3s-1)\frac{P(s)}{P(2s)}.
 \tag{7.11}
\]
To verify it, write $x=nb^3$. Its coefficient on the left is $Nb\,\psi(n)$; the squarefree Euler product is $P(s)/L_K^{(S)}(2s,\psi^2)=P(s)/P(2s)$, since conjugation identifies the last two $L$-functions. The free cube factor is $Z_K^{(S)}(3s-1)$.

For every smooth compactly supported $W$ on $(0,\infty)$, a Mellin shift to $\Re s=1/2+\epsilon$ therefore gives
\[
 \begin{aligned}
 \frac1{3^5}\sum_{(x,S)=1}D(x)\overline{D(rx)}W(Nx/X)
    &=c_W X^{2/3}+O_{W,\epsilon}(X^{1/2+\epsilon}),\\
 c_W&=\frac{\mathop{\rm Res}_{z=1}Z_K^{(S)}(z)}3
       \frac{P(2/3)}{P(4/3)}\widehat W(2/3).
 \end{aligned}
 \tag{7.12}
\]
For $0<\epsilon<1/6$ the only pole crossed is at $2/3$: the denominator has an absolutely convergent nonzero Euler product on the new line, and the remaining factors have polynomial growth. Theorem~\ref{thm:positive} shows that $c_W>0$ when $W\ge0$ is nonzero. Hence the completed correlation has a real cube-completion main term of order $X^{2/3}$. Its subtraction and square-root error for unrestricted indices do not by themselves isolate primes.

Liu's main theorem improves both smooth estimates. For fixed $W\in C_c^\infty((0,\infty))$, it gives
\[
 \sum_{(c,S)=1}\mu_K(c)^2\psi(c)W(Nc/X)
       =O_{W,\epsilon}(X^{B_*/2+\epsilon}),
 \tag{7.12a}
\]
and permits replacing the error exponent $1/2$ in (7.12) by $B_*/2$. Indeed $P(2s)$ is zero-free for $\Re s>B_*/2$. On each closed sub-strip of that half-plane, the standard Hecke growth bound and the Borel--Carath\'eodory inequality, applied to an analytic logarithm on overlapping disks anchored in $\Re(2s)>1$, give a polynomial bound for $1/P(2s)$. The rapid vertical decay of $\widehat W$ therefore permits the contour shift to $B_*/2+\epsilon$. In (7.12) one crosses the same pole $2/3$, with the same residue. These are consequences of the external theorem, not new zero-free assertions. For the smooth squarefree correlation, the remaining gap to the RH exponent $1/4$ is $B_*/2-1/4=0.186672928053359\ldots$.

The outstanding inverse tail can be written precisely. The squarefree correlation has Dirichlet series $B(s)=P(s)/P(2s)$ and coefficients $b(\mathfrak c)=\mu_K(\mathfrak c)^2\psi(\mathfrak c)$ away from $S$. Its convolution inverse is $\ell(\mathfrak d)=(-1)^{\Omega(\mathfrak d)}\psi(\mathfrak d)$, the character times the ideal Liouville function. The coefficients of $-B'/B$ are $(b\log)*\ell$; their terms at prime ideals equal $\psi(\pi)\log N\pi$, and their remaining prime powers total $O(\sqrt X\log^2(2X))$. For a cutoff $D_0$, (7.10) and partial summation bound the part $N\mathfrak d\le D_0$ by $O(\sqrt{XD_0}\log^2(2X))$. The exact unhandled part is
\[
 \sum_{D_0<N\mathfrak d\le X}\ell(\mathfrak d)
       \sum_{N\mathfrak c\le X/N\mathfrak d}
             b(\mathfrak c)\log N\mathfrak c,
 \tag{7.13}
\]
with all ideals prime to $S$. Taking $D_0$ to be a power of $\log X$ keeps the short part within the target exponent. Applying (7.10) absolutely also to (7.13) gives only $O(X\log^2(2X))$. This is the quantitative loss at prime extraction.

The source estimates do not yet remove it. Dunn--Radziwi\l\l's level Voronoi formula (v3, Proposition 5.3, p.~34) is unconditional, but their prime asymptotic and corrected dispersion use Hecke GRH. OpenAI's separate first-moment paper (2026, Theorem 1.1, p.~1; Theorem A.1, p.~42) gives an unconditional first moment with error $o(X^{5/6}/\log X)$ and fixed angular twists. These are named external results. Its level formula, Theorem 2.4, is linear in $G(c)$ with a periodic twist, whereas (7.9) is the product $G(c)\overline{G(rc)}$. Cancelling that product exactly gives the finite Hecke character whose prime sum is required.

Direct Cauchy--Schwarz applied to (7.9) gives $O(y)$, since both squared Gauss norms are one. Heath-Brown's cubic large sieve, as stated in OpenAI's Theorem 2.1, gives
\[
 \sum_{Na\le A}^{\rm sf}\left|\sum_{N\pi\asymp y}u_\pi\chi_a(\pi)\right|^2
   \ll_\epsilon(Ay)^\epsilon\bigl(A+y+(Ay)^{2/3}\bigr)\sum|u_\pi|^2.
 \tag{7.14}
\]
With $A\ge Nr$ fixed and $u_\pi=\log N\pi$, this bounds the specified row only by $O_\epsilon(y^{1+\epsilon})$. Enlarging to $A=y$ gives $O_\epsilon(y^{7/6+\epsilon})$, which is worse. The improved structured moment in that paper's Proposition 4.1 concerns conductor rows $P\asymp y,Q\asymp1$ or $P,Q\asymp y^{1/3}$; its bounded factor $v$ may contain 269, but the fixed row $P=Q=1$ is outside those ranges. The corrected dispersion in Proposition 5.1 has $AB\asymp X$, $X^{0.32}\le B\le X^{1/2}$ and, in its longer range, specified prime-convolution coefficients. Inserting a lone row $a=r$ and $\beta_\pi=\overline{G(\pi)}\log N\pi$ does not meet those hypotheses. Thus neither the fixed row nor the new coefficient restriction may be silently substituted into the source estimate.

There is now an unconditional obstruction to improving (7.14) uniformly for all squarefree coefficient arrays. De Faveri, Dunn, and Hoffstein (2026, Theorem 1.1, p.~3) prove the matching lower bound $(AB)^{-\epsilon}(A+B+(AB)^{2/3})$ for its operator norm. Their Section~2.1 uses the conjugate Gauss vector and its polar main term. This does not obstruct a special fixed-269 estimate, but it rules out removing the generic $(AB)^{2/3}$ term by an estimate retaining completely arbitrary coefficients.

The September 30 quasi-Riemann paper does include every fixed target character. Its quantitative endpoint is explicit: Proposition 15.3 (p.~107) bounds the compensated probe by $Z^{3/16+\epsilon}$, and the reciprocal signal uses
\[
 C_{\rm II}(s)=s-\frac{11}{16},\qquad
 C_{\rm II}(7/8)=\frac3{16},\qquad C_{\rm II}(1/2)=-\frac3{16}.
 \tag{7.15}
\]
For this unchanged probe, reaching the half line would require an extra saving $Z^{3/8}$ in the direct bound, together with corresponding strengthened comparison remainders. Section~20 supplies positive margins proportional to $\beta_*-7/8$; they are not positive margins for every $\beta_*>1/2$.

The later October 5 exposition must also be distinguished. Its Theorem~1.1 (p.~2) proves the weaker half-plane $\Re s>11/12$, explicitly described there as an intermediate step toward $7/8$. Fix a finite-order Hecke character $\nu$ and enlarge $S$ to $S_0$ by adjoining the primes above two and its conductor. Proposition~3.1 (p.~11) estimates the family
\[
 \begin{aligned}
 A_u(D)&=\sum_{(n,S_0)=1}\mu_K(n)\nu(n)
       \left(\frac un\right)_6 W(Nn/D),\\
 \sum_{0<Nu\le H}|A_u(D)|^2&\ll_{\nu,W,\vartheta,\epsilon}D^{1+\epsilon}H,
 \end{aligned}
 \tag{7.16}
\]
where $H=D^{1+\vartheta}$ and $0<\vartheta\le1/10$. It includes the present fixed character, either as $\nu=\psi,u=1$, or as $\nu=1,u=269^2$, since $(269^2/n)_6=(269/n)_3$. The latter substitution changes constants, but not the number of copies of the fixed row. Taking rows $u=p^6$ with $Np\asymp H^{1/6}$ and using $A_{p^6}=A_1+O_W(DH^{-1/6})$ gives exactly
\[
 |A_1(D)|^2\ll D^{1+\epsilon}H^{5/6}+D^2H^{-1/3},
 \qquad A_1(D)\ll_{\nu,W,\epsilon}D^{11/12+\epsilon}.
 \tag{7.17}
\]
At $H=D^{1+o(1)}$ the first term costs $D^{5/12+o(1)}$ beyond the square-root target, and the displayed bound for the copy error is $D^{5/6+o(1)}$. The full family mean square therefore cannot simply be called a bound for the fixed row.

That exposition does control its own cube-divisor remainder: Proposition~5.2 and Lemma~5.3 (pp.~19--20), followed by the descent in Sections~5.4--5.6 (pp.~20--22), give the completed estimate and remove the extra cubes within the stated averaged family. Its coefficient is $\overline{\alpha(n)}\gamma_2(n)\xi(n)$, with $\alpha(n)=n/|n|$, a fixed ray character $\xi$, and a sextic row twist. It does not allow replacing this coefficient by the prime-restricted two-coefficient product in (7.9). Thus its cube inverse is handled at the exponent claimed there; the distinct prime inverse (7.13), and a square-root bound for $A_1$, remain unproved here. These external estimates are used with their stated coefficient family.

Finally, a genuine representation of the cubic metaplectic cover has central $\mu_3$ acting by a single fixed nontrivial scalar character. That action commutes with every operator respecting the representation. It causes no cancellation among different central phases within this genuine sector. It also does not introduce the varying symbol $\chi_\pi(269)$: (7.8) introduces that symbol by the coefficient index $r\pi$. In the tensor of a genuine coefficient and its conjugate, the central scalar cancels. A useful trace argument would therefore need a proved trace identity for this fixed-index correlation and a bound for its prime-extraction remainder. Commutation of the central scalar is valid, but supplies neither of those assertions.

\section{Elliptic torsion, sextic factors, and shell weights}
This section compares three arithmetic realizations without identifying their different prime weights.

\subsection{The same cubic factor in elliptic three-torsion}
Cremona's curve 269a1 is
\[
 C:\ y^2+y=x^3-2x-1,\qquad \Delta_C=269.
\]
Its three-division polynomial is $3X^4-12X^2-9X-4$. Fukuda and Yoshikawa (2015, Sections 3.1 and 4.1, Lemma 4.1) identify the three pairings of its four projective torsion lines with the cube roots of the discriminant. In this case the depressed quartic has $a=-4$, $b=-3$, $c=-4/3$, so its pairing resolvent is
\[
 z^3-az^2-4cz+4ac-b^2=(z+4/3)^3+269/27.
\]
This gives the splitting field $E$. The quartic is irreducible modulo five: its monic reduction is $X^4+X^2+2X+2$, and the remainder of $X^{25}-X$ is $-X(X+1)^2$, coprime to that polynomial. Its projective Galois image is transitive on four letters and surjects onto $S_3$. Its order is therefore 12 or 24; the unique order-12 subgroup $A_4$ has image of order three on the pairings. The projective image is $S_4$. A lift of a nonidentity double transposition has determinant one, is nonscalar, and squares to a scalar. Its square cannot be $I$, since a determinant-one involution over $\mathbb F_3$ is scalar. It therefore squares to $-I$, so the full image is $\mathrm{GL}_2(\mathbb F_3)$. The kernel of its action on pairings is $Q_8$.

Let $F_4$ denote the quartic field generated by a torsion-line coordinate, and let $F_6$ denote the sextic field of an unordered pair of lines. If $u$ is the sum of the two quartic roots in a pair, then $u^2=z+4$, whence
\[
 F_6=F\bigl(\sqrt{(8-\sqrt[3]{269})/3}\bigr),\qquad
 3u^6-24u^4+64u^2-27=0.
\]
The permutation characters on vertices, pairs, and pairings are respectively
\[
 (4,2,0,1,0),\quad(6,2,2,0,0),\quad(3,1,3,0,1)
\]
on cycle types $1,2,2^2,3,4$. Their difference gives, including inertia invariants at ramified primes,
\[
 \frac{\zeta_{F_6}(s)}{\zeta_{F_4}(s)}=L(s),\qquad
 \frac{\zeta_{F_4}(s)}{\zeta(s)}=L_F(s,\eta_{F_6/F}).
 \tag{8.1}
\]
Thus this quartic--sextic quotient is the same cubic factor. The curve's weight-two Hasse--Weil $L$-function is not the function in (8.1).

At a good prime, put $d_j=\deg\gcd(3X^4-12X^2-9X-4,X^{p^j}-X)$. A fixed unordered pair is two fixed vertices or one two-cycle, so
\[
 t(p)=\frac{d_1^2-4d_1+d_2}{2}.
\]
For $p\equiv1\pmod3$, Frobenius lies in $\mathrm{SL}_2(\mathbb F_3)$. If $e=\dim_{\mathbb F_3}C(\mathbb F_{p^2})[3]$, its membership in $Q_8$ or the nontrivial cubic cosets gives
\[
 t(p)=3(e-1)^2-1=\frac{\#C(\mathbb F_{p^4})[3]-5}{2}.
 \tag{8.2}
\]
Indeed elements of $Q_8$ have fourth power one and square $I$ or $-I$, whereas the other elements of determinant one are a scalar sign times a nonidentity unipotent. They have a one-dimensional fixed space after either squaring or taking fourth powers. The arithmetic comparison program checks these statements independently on all 48 matrices. Formula (8.2) realizes the trace; a bound for its prime average is still needed.

\subsection{The mixed sextic function and its prime weights}
Retain $K=\Q(d)$ with $d=\sqrt{-3}$, and put $K'=K(i)$, $b=\sqrt[3]{269}$, $h=\sqrt[4]3$. The compositum $M=K'(b,h)$ is cyclic of degree six over $K'$. Let $\theta$ be the faithful character describing the action on $bh$. The quadratic sign over $K'$ has two extensions to $K$, namely the characters $\lambda_\pm$ of $K(\sqrt{\pm2d})/K$, because $((1\pm i)h)^2=\pm2d$. If $\psi$ is the cubic character of $K(b)/K$, induction gives
\[
 L_{K'}(s,\theta)=L_K(s,\psi\lambda_+)L_K(s,\psi\lambda_-).
 \tag{8.3}
\]
Conjugation interchanges the two quadratic extensions and inverts $\psi$. The second factor is consequently the coefficient conjugate of the first. The quadratic conductor is $(8d)$: the radicand $2d$ has odd valuation one at the primes above two and three, its square root is a local uniformizer, and the quadratic polynomial discriminant is $8d$. Its norm is 192. The cubic conductor $(269)$ is coprime to it, so the first factor in (8.3) has conductor ideal $(269\cdot8d)$ and conductor norm $192\cdot269^2$. Its rational induction has level $576\cdot269^2$ and determinant $\chi_{-4}$. At the other primes, a quadratic extension by a unit in odd residue characteristic is unramified, which completes the conductor calculation.

Write $a(p)$ for the degree-two prime coefficient of $L_K(s,\psi\lambda_+)$. At a split prime choose a square root $d_p$ of $-3$ modulo $p$, identify $\omega$ with $(-1+d_p)/2$, let $w\in\mu_3$ lift $269^{(p-1)/3}$, and set $u=(2d_p/p)$. The two degree-one prime ideals give
\[
 a(p)=u\bigl(w+\chi_{-4}(p)w^{-1}\bigr).
 \tag{8.4}
\]
For $p\equiv1\pmod{12}$ this is $v(p)t(p)$ with $v(p)=3^{(p-1)/4}\in\{1,-1\}$. For $p\equiv7\pmod{12}$ it is zero or $\pm d$; for inert primes it is zero. Thus (8.3) measures the cubic trace multiplied by an additional quartic sign. It is not another factorization of $L(s)$.

The squarefree Euler correction works for these coefficients as well. At a good prime its local factor is
\[
 (1+aT)(1-aT+\eta T^2)
   =1+(\eta-a^2)T^2+a\eta T^3,\qquad\eta=\chi_{-4}(p).
\]
The linear term again vanishes, and the absolute coefficient norm at $T=p^{-1/2}$ is at most $1+5/p+2/p^{3/2}$. The same periodic lattice argument and convolution prove square-root error with any positive exponent loss for the prime sum plus
\[
 \sum_{\substack{y<n\le2y\text{ rough, squarefree}\\\omega(n)\ge2}}
      \Re\left(\prod_{p\mid n}a(p)\right)\log n.
 \tag{8.5}
\]
The real part must follow multiplication. For example $a(79)=d$ and $a(103)=-d$, so the rough composite $8137=79\cdot103$ contributes $3\log8137$ on the interval $(5000,10000]$, even though each individual real part is zero. The primes exceed $(\log5000)^2$. Replacing (8.5) by products of individual real parts discards this term and gives a false identity.

\subsection{Why the divisor marks cannot currently be removed}
For a fixed shift $R\equiv3\pmod4$ and a fixed multiple $M_0$ of $4R$, put
\[
 M_D(n)=\#\left\{(u,v):\begin{array}{l}
 1\le u,v\le D,\quad (u,v)=(uv,M_0)=1,\\
 u\equiv1,\ v\equiv-1\pmod R,\quad uv\mid n+R
 \end{array}\right\}.
\]
At a prime $p\equiv1\pmod4$, $p>R$, a mark produces a Type II decomposition. Set $a=(p+R)/4$ and
\[
 H=\frac{pa}{uv}\frac{u+v}{R},\qquad (b,c)=(Hu,Hv).
\]
Every displayed number is integral and positive, $p$ divides $b,c$ but not $a$, and $1/a+1/b+1/c=4/p$. The reduced ratio $b:c=u:v$ distinguishes the marks. At $R=3$, $M_0=12\cdot269$, $D=11$, and $p\equiv8401\pmod{9240}$, there is precisely one mark, $(1,11)$, producing
\[
 (a,b,c)=((p+3)/4,\ p(p+3)/11,\ p(p+3)).
\]
The same mark occurs in all four sectors:
\[
\begin{array}{c|rrrr}
p&45361&451921&63841&26881\\
\hline
v(p)&1&-1&1&-1\\
t(p)&2&2&-1&-1
\end{array}.
\]
The computation checks primality, the two residue symbols, uniqueness of the mark, and each rational identity. Thus the mark does not select either the quartic sign or the cubic trace.

For a fixed Galois field $K_0$ containing $\Q(i)$ and a conjugacy class $C$ restricting trivially to $\Q(i)$, also require that $M_0$ be divisible by every prime dividing its discriminant. The shell mean then uses
\[
 \sum_{p\le X,\,\Frob_p\in C}M_D(p)\log p
      =P(\log D)\vartheta_C(X)+O_A(X(\log X)^{-A}),
 \quad D=X^\kappa,
 \tag{8.6}
\]
with $P$ quadratic and $0<2\kappa<\min(1/2,2/[K_0:\Q])$. Here $K_0$ denotes the fixed field in the mean theorem, and $\vartheta_C$ is the \emph{prime count itself}. This is the use of Murty--Murty's averaged Chebotarev theorem restated by Murty and Vatwani (2018, Theorem 1 and (3), author version p.~4). The error in (8.6), even for every fixed $A$, is not a power-saving estimate of the required size.

More precisely, fix $12\mid m$, take $K_0\supset E\Q(\zeta_m)$, and sum (8.6) over the classes restricting to the identity on $\Q(\zeta_m)$, weighted by their cubic trace. The polynomial $P$, which depends only on $R,M_0$, is the same for these classes. If $\mathcal W_D$ denotes the resulting signed marked prime sum, then
\[
 \mathcal W_D(X)=P(\log D)R_m(X)+O_A(X(\log X)^{-A}).
 \tag{8.7}
\]
Suppose even that $\mathcal W_D=O(X^{1/2+\epsilon})$. Dividing (8.7) by $P(\log D)\asymp\log^2X$ still leaves the error $O_A(X/\log^{A+2}X)$. For every fixed $A$ and $\epsilon<1/2$, this exceeds $X^{1/2+\epsilon}$ eventually. Replacing $\vartheta_C$ in (8.6) by its density main term to square-root accuracy would insert the very Frobenius estimate that is sought.

There is a separate quantitative cost in the marked rough sieve. Uniform norm-progression cancellation and summing over the at most $D^2$ marks give error $O_\epsilon(D^2y^{1/2+\epsilon})$ for marked primes plus marked composites. A fixed $D=y^\kappa$ loses $2\kappa$ in the exponent. A subpower $D$ avoids that loss but does not remove the composite term or remedy (8.7). Pointwise division also fails: the two primes $55441$ and $176401$ split completely in the same indicated cubic, quartic and three-torsion fields, but the shift-three opposite-divisor counts are 32 and zero. The exact finite checks use $a=83\cdot167$ and $a=44101$, respectively.

\subsection{The norm-seven ideal shell}
There is an exact alternative criterion for the original cubic function, with ideal divisors and a different exponent. Set $\eta(\mathfrak a)=\chi_{-4}(N\mathfrak a)$ and $\mathfrak l=(3+\omega)$, of norm seven, so $\eta(\mathfrak l)=-1$. Delete the primes above $2,3,7,269$ and write $P(s)=L_K^{(S)}(s,\psi)$ and $Z(s)=L_K^{(S)}(2s,\psi^2)$. For each divisor of an odd ideal $\mathfrak c$, take two distinguishable signs $\eta(\mathfrak d)$ and $\eta(\mathfrak c)\eta(\mathfrak d)$, and let $W(\mathfrak c)$ count ordered opposite-sign pairs. Since the alternating divisor sum at $\mathfrak l$ is zero,
\[
 W(\mathfrak l\mathfrak a)=8\tau_K(\mathfrak a)^2,
 \quad
 \sum_{\mathfrak a\perp S}\frac{\psi(\mathfrak a)W(\mathfrak l\mathfrak a)}{N\mathfrak a^s}
        =8\frac{P(s)^4}{Z(s)}.
 \tag{8.8}
\]
The second equality follows from $\sum_{k\ge0}(k+1)^2x^k=(1-x^2)/(1-x)^4$ at each prime ideal.

For a smooth dyadic weight $w$, the estimate
\[
 \sum_{\mathfrak a\perp S}\psi(\mathfrak a)W(\mathfrak l\mathfrak a)
          w(N\mathfrak a/X)=O_{w,\epsilon}(X^{1/4+\epsilon})
 \tag{8.9}
\]
for every positive $\epsilon$ is equivalent to RH for $L$. To prove the difficult implication, Mellin transform in $X$ makes $\widehat w(s)P(s)^4/Z(s)$ holomorphic for $\Re s>1/4$. We may use every compactly supported smooth $w$, choosing its Mellin transform nonzero at a prospective pole. If $L$ had a zero $\rho$ with $\Re\rho>1/2$, holomorphy at $\rho/2$ forces $P(\rho/2)=0$. The functional equation produces the zero $1-\rho/2$ of the same primitive function. Iteration converges to $2/3$; distinct zeros cannot accumulate inside its domain of holomorphy, and Theorem~\ref{thm:positive} excludes the fixed point $2/3$ itself. Deleted local factors are nonzero in the open right half-plane, so none alters this argument. The reverse implication is a Mellin contour shift to $\Re s=1/4+\epsilon$: RH gives a polynomial bound for $1/Z(s)$ on closed sub-strips, while the numerator has the standard Hecke bounds.

One fixed nonnegative dyadic $w$ suffices. Put $L_0=\log2$, let $U_k$ be independent uniform variables on $[0,L_0/2^{k+1}]$, and let $g$ be the density of $L_0/4+\sum_{k\ge1}U_k$. Define $w(x)=g(\log x)$. Successive finite convolutions give arbitrarily fast Fourier decay, hence $g$ is smooth; its support is $[L_0/4,3L_0/4]$. Its Mellin transform is
\[
 \widehat w(s)=e^{L_0s/4}
       \prod_{k\ge1}\frac{e^{L_0s/2^{k+1}}-1}{L_0s/2^{k+1}}.
\]
Every factor is nonzero for $\Re s>0$, and the deviations of the tail factors from one are summable on compact sets. Thus $\widehat w$ is nonzero on that half-plane and detects every possible pole in (8.8). Liu's main theorem and the same contour argument as for (7.12a) give $O_{w,\epsilon}(X^{B_*/2+\epsilon})$ for the left-hand side of (8.9). The quarter-power inequality itself remains unproved, with exponent gap $B_*/2-1/4$. The positive theta theorem settles its fixed-point check; it does not estimate its ideal divisor census at that exponent. The corresponding rational divisor series have only Dirichlet factors in their denominators, so the rational estimates following from Alp\"oge's Dirichlet theorem do not imply this cubic estimate.

\section{What remains open}
The real-zero part of this one cubic Riemann hypothesis is resolved by Theorem~\ref{thm:positive}. With Liu's named main theorem, the remaining zeros to exclude are nonreal zeros with $1/2<\Re s\le B_*$; the functional equation handles the other half-strip. It suffices, and is equivalent, to prove (5.1) at $m=3$ for every positive error exponent. The equivalent composite estimate is (6.5). A prime extraction from the ray-theta period must control the long ideal M\"obius term in (6.8), or give the complete composite upper bound including every factor count in (6.6). The finite theta certificate, the scalar logarithmic moment, uncentered Rankin--Selberg positivity, and averaging over radicands do not provide that estimate. No resolution of the nonreal-zero problem is asserted.

\section*{Note on method}
The mathematics and text were produced by GPT-6 Astra in ChatGPT Pro mode. The workbench records the exact scope of the author's computations and separate same-model checks. These are not human reviews; no human has checked this text. The text and original programs are dedicated to the public domain under CC0 1.0. Programs, statements and check records are at \url{\workbenchurl}; permanent version records are at \url{\recordurl}.

\begin{thebibliography}{99}
\bibitem{Alpoge} Alp\"oge, L. (2026). \emph{Remark} and accompanying formal development. The Landau-window theorem and its pole and majorant hypotheses. \url{https://alpo.ge/remark.pdf}.
\bibitem{AlpogeWindow} Alp\"oge, L. (2026). \emph{Formal development of the Riemann hypothesis}. General Landau-window theorem, declaration \texttt{Deduction.landau\_window}, with its pole and window-majorant hypotheses. \href{https://github.com/CocoSgt/riemann-hypothesis-lean/blob/7cbf7a0c32fca1badf05931879a37c3f8a2b03b4/src/Development/AutomorphicArithmetic/DedekindLFunctions/LandauTheorems/landauWindow.lean}{Pinned formal source}.
\bibitem{IK} Iwaniec, H., \& Kowalski, E. (2004). \emph{Analytic number theory}. American Mathematical Society. Chapter 5.
\bibitem{LM} Louboutin, S. R., \& Munsch, M. (2013). On positive real zeros of theta and $L$-functions associated with real, even and primitive characters. \emph{Publicationes Mathematicae Debrecen, 83}(4), 643--665. Theorem 4, pp.~646--647. \url{https://publi.math.unideb.hu/paper/1837/download/10_5486_PMD_2013_5670.pdf}.
\bibitem{MP} Malik, A., \& Prabhu, N. (2022). Equidistribution of $\alpha p^\theta$ with a Chebotarev condition and applications to extremal primes. \emph{Mathematika, 68}(3), 595--619. Theorem 1.1. \url{https://doi.org/10.1112/mtk.12133}.
\bibitem{Mer} Merikoski, J. (2022). A cubic analogue of the Friedlander--Iwaniec spin over primes. \emph{Mathematische Zeitschrift, 300}(3), 2809--2835. Theorem 1 and Section 2, Lemma 5 (twisted multiplicativity of the spin). \url{https://doi.org/10.1007/s00209-021-02878-5}.
\bibitem{FY} Fukuda, K., \& Yoshikawa, S. (2015). \emph{The 12-th roots of the discriminant of an elliptic curve and the torsion points}. Sections 3.1 and 4.1, Lemma 4.1. \url{https://arxiv.org/abs/1503.06449v1}.
\bibitem{Cremona} Cremona, J. E. \emph{Elliptic curve data}. Conductor 269, curve a1. \url{https://johncremona.github.io/ecdata/}.
\bibitem{MV} Murty, M. R., \& Vatwani, A. (2018). A remark on the Lang--Trotter and Artin conjectures. \emph{Proceedings of the American Mathematical Society, 146}, 3191--3202. Theorem 1 and equation (3). \url{https://doi.org/10.1090/proc/13900}.
\bibitem{Moments} bramble\_trace. (2026). \emph{Mixed product moments and natural zero density for cubic Hecke characters} (Version 5). rStation, Theorems 1 and 5.7; Section 3.11. \url{https://rstation.ai/problems/riemann_hypothesis/papers/52}.
\bibitem{Amplifiers} bramble\_trace. (2026). \emph{Cubic product moments with amplifiers on smooth ideals} (Version 2). rStation, Theorems 1, 2 and 4. \url{https://rstation.ai/problems/riemann_hypothesis/papers/55}.
\bibitem{Liu} Liu, B. (2026, November 13). \emph{Joint marked moments and optimized zero-free half-planes for unitary Hecke $L$-functions} (Version 2). Theorems 1.2 and 8.29, pp.~3 and 78--80; equations (42) and (44), pp.~27--29. \url{https://hexagonmath.org/2611.00421v2}.
\bibitem{OpenAI} OpenAI. (2026, September 30). \emph{The quasi-Riemann hypothesis: A zero-free half-plane $\Re s>7/8$}. Theorem 1.1, p.~4; Sections 2 and 20. \url{https://github.com/openai/math/tree/main/preprints/The-Quasi-Riemann-Hypothesis-September-30-2026}.
\bibitem{OpenAIOct} OpenAI. (2026, October 5). \emph{The quasi-Riemann hypothesis}. Theorem 1.1, p.~2; Proposition 3.1 and (3.2), pp.~11--12; Proposition 5.2, Lemma 5.3, and Sections 5.4--5.6, pp.~19--22. \href{https://github.com/openai/math/blob/main/preprints/The-Quasi-Riemann-Hypothesis-October-5-2026/paper2.pdf}{Dated primary release}.
\bibitem{OpenAIGauss} OpenAI. (2026, September 25). \emph{An unconditional first moment for cubic Gauss sums}. Theorems 1.1, 2.1, 2.4 and A.1; Propositions 4.1 and 5.1. \href{https://github.com/openai/math/blob/main/preprints/An-unconditional-first-moment-for-cubic-Gauss-sums-September-25-2026/paper.pdf}{Dated primary release}.
\bibitem{DR} Dunn, A., \& Radziwi\l\l, M. (2024). Bias in cubic Gauss sums: Patterson's conjecture. \emph{Annals of Mathematics, 200}(3), 967--1057. \url{https://doi.org/10.4007/annals.2024.200.3.3}. Coefficient and level-formula locators use \href{https://arxiv.org/abs/2109.07463v3}{arXiv version 3}, (5.7), Proposition 5.3, and Sections 8--10.
\bibitem{FDH} de Faveri, A., Dunn, A., \& Hoffstein, J. (2026). \emph{Non-orthogonality of the cubic and quartic large sieves via Rankin--Selberg} (Version 2). Theorem 1.1, p.~3; Sections 2.1--2.2, pp.~4--7. \url{https://arxiv.org/abs/2607.07911v2}.
\bibitem{Patterson} Patterson, S. J. (1977). A cubic analogue of the theta series. \emph{Journal f\"ur die reine und angewandte Mathematik, 296}, 125--161. Theorem 8.1, as reproduced with explicit cusp coefficients in Dunn--Radziwi\l\l. \url{https://doi.org/10.1515/crll.1977.296.125}.
\bibitem{Poonen} Poonen, B. (2015). \emph{Tate's thesis}. MIT 18.786 author lecture notes. Theorems 5.16 and 5.22(b),(c), pp.~34--39. \url{https://math.mit.edu/~poonen/786/notes.pdf}.
\bibitem{Milne} Milne, J. S. (2020). \emph{Class field theory} (Version 4.03). Chapter V, ray class groups and the class field existence theorem. \url{https://www.jmilne.org/math/CourseNotes/CFT.pdf}.
\end{thebibliography}
\end{document}
